\ifdefined\MyWebBook\else
\documentclass[../textbook.tex]{subfiles}
\begin{document}
\fi
\bookchapter{上下文学习}{ch:context-learning}

一个固定参数的语言模型，在收到不同指令、示例和背景信息时，会给出不同的条件预测。增加几个示例、要求分解任务，或生成多个候选再选择，都会改变求解过程；这些策略分别作用于输入信息、生成路径和决策规则，其质量收益应与额外计算成本共同评价。

比较这些策略时，应固定任务集和计算预算，并记录提示模板、示例、采样设置与答案提取规则，以便复现结果。

\section{固定模型上的条件学习}

\subsection{输入条件与参数更新}
\term{上下文学习}{In-Context Learning}{ICL}通过输入中的任务描述与示例改变模型的预测，当前请求不更新参数\cite{brown2020fewshot}。记参数为 $\theta$，当前输入为 $x$，有序示例序列为 $D=((x_1,y_1),\ldots,(x_m,y_m))$，其预测写为
\begin{equation}
p_\theta(y\mid x,D,I,\mathcal B,O,T),
\label{eq:icl-condition}
\end{equation}
其中 $I$ 为任务指令，$\mathcal B$ 为背景，$O$ 为输出约束，$T$ 为消息序列化与生成起始协议。自然语言文字不变时，角色分隔符和模板改变仍会改变输入模型的词元序列。

\term{零样本提示}{Zero-Shot Prompting}{}没有当前任务的示例对，但可以包含指令和背景。\term{少样本提示}{Few-Shot Prompting}{}额外提供少量示例；样本数指当前上下文内的示例数量。ICL 的表现可能来自对已学任务的识别、格式适应、信息提取或新映射的临时推断；成功的回答与上述机制之间没有一一对应。

\bookxref{ch:parameter-efficient-finetuning}中的连续提示训练需要更新任务参数，而上下文学习中的文本示例\emph{完全不改变底层模型权重}。\emph{背景知识}提供外部事实依据，而\emph{示范示例}展示输入到输出的映射结构；两者在语义职责上严格区分。两者可以同时存在，来源、正确性与评价方式需分别记录。

\begin{table}[htbp]
\centering
\caption{本章主要符号与对象}
\label{tab:context-learning-symbols}
\begin{tabular}{ll}
\toprule
符号 & 含义\\
\midrule
$\theta,x,D$ & 固定模型参数、任务输入、有序示例\\
$C$ & 完整运行时上下文，包括模板与生成起点\\
$z,a$ & 中间步骤文本与最终答案\\
$q(a,z\mid C)$ & 实际解码策略产生的联合分布\\
$h(z,a)$ & 版本化答案解析与规范化函数\\
$K,n_a$ & 候选数量与答案 $a$ 的票数\\
$v(C,z,a)$ & 验证器对候选的评分\\
$L_C,L_i,B$ & 上下文长度、候选输出长度、预算\\
\bottomrule
\end{tabular}
\end{table}

\subsection{上下文结构与组件职能}
任务指令说明要执行的操作、判断标准与适用边界；示例具体呈现输入输出关系；背景提供任务所需信息；输出约束规定返回内容及可解析形式。它们共同形成 $C$，四者在权威层级上各不相同。背景中的引用、检索结果与用户提供的待处理文本应作为数据解释，命令句式出现在数据位置也不改变它的身份。

明确边界既防止指令冲突，也减少统计歧义。例如分类任务的标签应是固定集合，同一标签在示例中对应不同含义会直接造成矛盾；结构化输出应规定字段类型、单位和缺失值。格式检查确定结果能否解析，事实核验与业务规则检查确定内容是否有效。

\begin{figure}[htbp]
\centering
\begin{tikzpicture}[>=stealth,every node/.style={font=\small},b/.style={draw,align=center,minimum width=2.35cm,minimum height=.85cm}]
\node[b](i)at(0,1.2){任务指令};\node[b](d)at(0,-.2){有序示例};\node[b](b)at(3.2,1.2){任务背景};\node[b](o)at(3.2,-.2){输出约束};
\node[b](c)at(6.3,.5){模板序列化\\上下文 $C$};\node[b](m)at(9.5,.5){固定模型\\条件预测};
\draw[->](i)--(c);\draw[->](d)--(c);\draw[->](b)--(c);\draw[->](o)--(c);\draw[->](c)--(m);
\end{tikzpicture}
\caption[上下文学习输入组成与预测过程]{上下文学习输入组成与预测过程。不同上下文组件进入同一条件预测过程；当前请求不改变模型参数。}\label{fig:icl-context}
\end{figure}

\section{上下文样本配置}

\subsection{相关性、覆盖和预算}
\term{示例库}{Example Bank}{}是可复用的示范资产，应记录来源、任务类型、答案、用途和版本。候选首先满足访问权限与数据质量，再进行相关性排序。把无权使用的信息放入模型后再要求其“不要泄露”，其效果取决于模型是否遵守，前置隔离才是可靠手段。

示例相关性可以依据词汇重合、语义向量、任务结构或已知难度。示例检索不仅应考量词表或表层嵌入的重合度，更须对齐底层解题逻辑与约束结构：两道题目可能仅替换实体名称，却需要相反的判定规则；表层形式差异显著的问题也可能共享相同的约束推理拓扑。因此，示例选择准则应联合编码语义相似度与任务求解结构。

一种显式的设计目标是，在候选集合 $\mathcal E$ 中选择子集 $S$，最大化
\begin{equation}
F(S;x)=\sum_{e\in S}r(e,x)+\lambda\operatorname{Coverage}(S)
-\mu\sum_{\{e,e'\}\subseteq S}\operatorname{Redundancy}(e,e'),
\quad \sum_{e\in S}\ell_e\le B_D.
\label{eq:icl-selection}
\end{equation}
$r$ 衡量相关性，覆盖项表示不同必要技能或标签，冗余项惩罚重复，$\ell_e$ 为序列化后的示例长度。它是用于明确取舍的设计形式，任务最优性需要另行验证。若所有候选相关性均不足，可回退到更少示例或零样本，为凑够固定数量硬塞无关样本只会增加噪声。

例如预算 $B_D=200$ 个词元，三个示例长度为 $(120,100,80)$，相关性为 $(0.9,0.8,0.6)$。仅按相关性从高到低且放不下就跳过，会选第一、第三条，总相关性 $1.5$；第二、第三条仅需 $180$ 个词元、总相关性 $1.4$，若二者覆盖两个互补技能而第一、第三条重复，覆盖项可以使其更优。示例选择使用预先规定的信息与评分规则。

\subsection{上下文顺序效应}
模型接收的是有序序列，所以
\begin{equation}
p_\theta(y\mid x,D)=p_\theta(y\mid x,\pi(D))
\end{equation}
一般只对特定排列 $\pi$ 成立，而不对任意排列成立。位置编码、因果注意力与前文示例形成的局部模式使次序具有作用。已有研究直接观察到少样本顺序敏感性\cite{lu2022orderedprompts}；但某个模型上的最好排列在其他模型或任务上需要重新测量。

长度处理会改变模型实际接收的示例。若输出预算固定，增加示例可能挤掉必要背景；若上下文被截断，也可能留下不完整的示例或答案。因此，复现条件输入时应保存截断后的词元序列。\footnote{示例前缀可被缓存以降低重复预填充成本，但缓存改变的是执行复用，不会消除示例在上下文窗口中占用的位置。}

\subsection{标签分布和内容先验}
示例同时展示格式、输入分布、标签空间和具体映射。针对部分分类与多项选择任务的研究表明，这些因素可以分别影响 ICL 结果\cite{min2022demonstrations}；示例标签的作用随任务而变化。若任务需要通过示例学习陌生符号映射，错误标签可能直接破坏映射识别。

\term{标签词}{Verbalizer}{}把类别对应到模型输出文本。不同标签词的预训练频率、词元长度和上下文位置会引入偏好。若标签含多个词元，应该比较完整标签序列的概率及必要结束边界，首词元概率只覆盖其中一段。用平均词元对数概率替代序列概率又会改变决策目标，需明确声明。

一种\term{上下文校准}{Contextual Calibration}{}方法，用无任务内容的输入估计标签偏好 $b_a>0$，再对目标输入概率作比值校正\cite{zhao2021contextcalibration}：
\begin{equation}
\widetilde p(a\mid x,C)=\frac{\frac{p_\theta(a\mid x,C)}{b_a}}{\sum_{a'}\frac{p_\theta(a'\mid x,C)}{b_{a'}}}.
\end{equation}
例如两标签的原概率为 $(0.6,0.4)$、偏好估计为 $(0.8,0.2)$，比值为 $(0.75,2)$，归一化后为 $(\frac{3}{11},\frac{8}{11})$。校正使标签排序反转。无内容输入是否能代表应扣除的偏好、是否需要平滑，以及真实类别先验是否本来就不均衡，都会影响适用性。校准系数应在开发数据上确定。

\begin{algorithm}[预算内的示例构造]
\AlgoInput{问题、版本化示例库、选择策略、上下文上限与保留输出长度。}
\AlgoOutput{实际序列化上下文及示例清单。}
\AlgoState{可用词元预算、已选集合、覆盖状态。}
\begin{enumerate}
\item 过滤不符合权限、任务与质量条件的候选；排除与保留评估样本不允许的家族重叠。
\item 扣除指令、问题、必要背景、模板和输出预留，得到示例预算；预算不足时采用明确的缩减或失败策略。
\item 按预定相关性、覆盖和冗余准则选择可容纳示例，采用固定平局规则。达到预算或无合格候选时停止。
\item 按已声明顺序构造完整消息并实际分词；如长度变化导致超限，整条移除最低优先级示例后重新构造，不截断到示例答案中间。
\end{enumerate}
\textbf{不变量：}最终上下文长度符合预算，示例来源与顺序可追溯，选择过程没有使用当前测试题的标准答案。
\end{algorithm}

\section{分步生成的概率结构}

\subsection{中间推理文本的作用}
\term{思维链提示}{Chain-of-Thought Prompting}{CoT Prompting}通过步骤示范或任务分解要求，引导模型先生成中间文本再形成答案\cite{wei2022cot}。令中间文本为 $z$、答案为 $a$，在同一生成协议下有
\begin{equation}\label{eq:icl-marginal}
 \begin{aligned}
  p_\theta(z,a\mid C) &= p_\theta(z\mid C)p_\theta(a\mid C,z), \\
  p_\theta(a\mid C) &= \sum_z p_\theta(z,a\mid C).
 \end{aligned}
\end{equation}
中间词元成为后续位置的条件，能够保存局部结果、展开约束和组织计算。中间文本还可能包含错误、遗漏或事后解释，其有效性通过步骤与结果检查。

式\eqref{eq:icl-marginal}描述允许生成 $z$ 的协议。要求直接回答可能改变提示和输出语法，其答案分布不必等于先生成步骤再边缘化的分布。比较两种协议时，应同时记录提示、输出语法与答案分布。

单轨迹生成访问一个候选 $z$；边缘概率则汇总所有路径。重复前提、同义复述或复制错误会耗费预算却不改善答案。对于可外部执行的算术、程序和约束，应把验证对象写成明确输入、操作与结果，据此检查步骤有效性。

\subsection{路径最大概率与答案最大概率}
设四条完整路径的概率分别为 $0.30,0.25,0.25,0.20$。第一条给出答案 B，第二、第三条给出答案 A，第四条给出答案 C。最大概率路径给出 B，但答案边缘概率为 $p(A)=0.50,p(B)=0.30,p(C)=0.20$，最大概率答案是 A。

\keyconcept{单路径最优搜索与边缘化答案聚合对应不同的概率目标：精确搜索联合概率最高的单条推理路径，并不保证收敛至边缘概率最高的终态答案。}单条解答的最优搜索与边缘化答案空间的众数聚合对应不同的概率极值，而逐词贪心解码亦无法保证得到全局概率最优的完整轨迹；相关搜索与剪枝边界参见\bookxref{ch:autoregressive}。

\section{自一致性解码}

\subsection{采样分布估计}
\term{自一致性}{Self-Consistency}{}生成多条候选路径，提取最终答案并按出现频率聚合\cite{wang2022selfconsistency}。令真实解码分布为 $q(z,a\mid C)$，答案解析函数为 $h$，则 $K$ 次采样的计数为
\begin{equation}
 \begin{aligned}
  n_b &= \sum_{i=1}^{K}\mathbf1[h(z_i,a_i)=b], \\
  \widehat q_K(b) &= \frac{n_b}{K}.
 \end{aligned}
\end{equation}
在独立同分布采样下，$\mathbb E[\widehat q_K(b)]=q_h(b)$，方差为 $\frac{q_h(b)(1-q_h(b))}{K}$，其中 $q_h$ 是经过解析映射后的答案分布。温度、top-$k$、top-$p$ 或停止规则改变 $q$，因此投票估计的是经解码规则和答案解析映射得到的分布，通常不同于原始模型分布 $p_\theta$。

重复执行确定性贪心通常只得到同一个答案，候选数量在这里不提供新的采样信息。反之，同一个固定模型使用独立随机数，可以在给定问题与提示的条件下产生独立样本；样本独立性取决于条件分布与随机采样过程。跨题的共同偏差、候选复用同一已生成前缀、相互参考修改或共享随机状态，则会改变应采用的独立性假设。

\begin{assumption}[二元多数正确率]
固定一题，候选的正确性是独立同分布 Bernoulli 变量，单次正确率为 $p$。候选数 $K$ 为奇数，决策仅区分正确和错误两类；不存在解析失败、拒答或多种错误标签竞争。
\end{assumption}

满足上述假设时，多数正确概率为
\begin{equation}
P_K=\sum_{j=\frac{K+1}{2}}^{K}\binom Kj p^j(1-p)^{K-j}.
\label{eq:icl-majority}
\end{equation}
$p=0.6,K=5$ 时，$P_5=10(0.6)^3(0.4)^2+5(0.6)^4(0.4)+(0.6)^5=0.68256$。$p>\frac{1}{2}$ 是这一二元改善机制的关键；$p<\frac{1}{2}$ 时，增加候选会更稳定地趋向错误答案。

真实答案空间一般不是二元。若正确答案概率 $0.4$，其他答案分别为 $0.35$ 和 $0.25$，正确答案虽然不足半数，却是唯一众数；独立大量采样的最高票可能收敛到正确答案。若一个错误答案拥有最大概率，则自一致性趋向于强化它。开放答案空间的投票结果取决于各答案的概率与候选采样方式。

\subsection{相关性、平票与解析失败}
若正确性指示变量具有相同边际概率 $p$、任意两个的相关系数为 $\rho$，则平均正确率方差为
\begin{equation}
\operatorname{Var}(\bar C)=\frac{p(1-p)}K[1+(K-1)\rho].
\end{equation}
该式只在所声明的等相关结构下成立，而且给出的是平均正确率的方差而非相关投票的完整尾概率；仅知道 $\rho$ 还不足以代入式\eqref{eq:icl-majority}。候选共享一种错误解释时，经验收益可能很早饱和。

\term{答案规范化}{Answer Canonicalization}{}把单位、数值格式和等价表达映射到任务定义的同一答案。例如在要求以米给出长度的题中，$1.5$ 米与 $150$ 厘米可经单位换算合并；不带单位的 $150$ 缺少量纲依据，需要单独判定。大小写、空白、小数容差与集合顺序是否可忽略，取决于任务本身。

\term{解析失败}{Parse Failure}{}应作为显式状态并保留在分母中。若把失败结果记为 $\bot$，则所有计数之和仍为 $K$。可以只在有效答案中选择，但必须同时报告有效覆盖率 $1-\frac{n_\bot}{K}$、总候选上的支持率和有效候选中的支持率。\footnote{把无效输出排除后计算的正确率是条件于“输出有效”的指标；它用于诊断解析器，覆盖全部请求的端到端正确率需要另行报告。}

奇数候选仍可能出现多类别平票。五个候选可以形成 $2:2:1$。预先定义的平票处理、最低有效数与最低票差可以触发继续采样、验证或拒答；事后挑选更符合标准答案的候选会使规则失去意义。票数领先描述样本共识，正确概率需在独立数据上校准。

\section{候选答案验证}

\subsection{投票排序和验证排序}
\term{验证器}{Verifier}{}检查候选的某些性质或为其正确性评分。\term{多候选择优}{Best-of-$K$ Selection}{Best-of-$K$}可写为
\begin{equation}
 \begin{aligned}
  i^* &= \arg\max_{1\le i\le K}v(C,z_i,a_i), \\
  \widehat a &= h(z_{i^*},a_{i^*}).
 \end{aligned}
\end{equation}
这与票数选择不同：罕见答案可能通过严格验证，而最高票答案可能违反约束。训练验证器从多个数学候选中选择，是一条已有研究路线\cite{cobbe2021verifiers}；选择器的收益取决于验证器的误接受率与误拒绝率。

若每次独立生成正确候选的概率为 $p$，则至少有一个正确候选的概率为
\begin{equation}
P(\text{覆盖正确答案})=1-(1-p)^K.
\label{eq:icl-coverage}
\end{equation}
$p=0.6,K=5$ 时为 $0.98976$，远高于前面的多数正确率，但这是完美选择器能够达到的候选覆盖，实际系统正确率还要打上验证器误差的折扣。正确候选已经出现时，验证器仍可能选错。扩大 $K$ 还可能引入更多能欺骗评分器的高分错误，固定验证器的选择质量因此随 $K$ 非单调变化。

\subsection{验证器适用范围}
形式验证、程序执行、量纲检查、数据库约束和外部证据分别检查不同性质。算术验证同时检查表达式与题意的对应关系，程序测试覆盖指定输入，引用核验检查资料对主张的支持。验证报告记录对象、方法与覆盖范围。

模型自评受限于基座的内在知识与诱导偏差：模型回读自身输出虽能修正局部格式，但极易沿用生成时的系统性误解。若评估模型与生成模型共享同一组参数或训练分布，表面上的两阶段调用在概率推断上难以提供正交的外部验证支撑。\term{拒答}{Abstention}{}应作为系统的合法终态，用于覆盖共识不足、验证条件缺失或关键证据不可得的场景；强制返回置信度不足的最高分候选将掩盖真实的不确定性风险。

\begin{figure}[htbp]
\centering
\begin{tikzpicture}[>=stealth,every node/.style={font=\small},b/.style={draw,align=center,minimum width=2.4cm,minimum height=.85cm}]
\node[b](g)at(0,0){多条候选};\node[b](p)at(3.1,0){解析与规范化};
\node[b](vote)at(6.5,1.2){自一致性\\按答案计数};\node[b](verify)at(6.5,-1.2){验证选择\\按证据或分数};\node[b](r)at(10,0){决策或拒答};
\draw[->](g)--(p);\draw[->](p)--(vote);\draw[->](p)--(verify);\draw[->](vote)--(r);\draw[->](verify)--(r);
\end{tikzpicture}
\caption[聚合与验证决策分支对比]{聚合与验证决策分支对比。聚合与验证是不同决策分支，可以组合，但应分别记录各自依据。}\label{fig:icl-selection-flow}
\end{figure}

表\ref{tab:prompt-taxonomy}比较了提示工程与测试时计算（Test-Time Compute）中的几种方法。它们采用不同的候选组织和选择方式，适用性取决于任务是否便于分解、候选能否验证，以及可用的计算预算。

\begin{table}[htbp]
\centering
\footnotesize
\setlength{\tabcolsep}{3pt}
\caption{提示与测试时计算方法对比}
\label{tab:prompt-taxonomy}
\begin{tabularx}{\linewidth}{@{}p{20mm}p{16mm}cp{28mm}X@{}}
\toprule
技术范式 & 搜索拓扑 & 相对计算 & 核心机制与决策依据 & 典型适用场景 \\
\midrule
Zero / Few-Shot & 单步映射 & $1\times$ & 先验与少样本上下文对齐 & 分类、抽取与常规问答 \\
思维链 (CoT) & 线性单链 & $1\times$ (长输出) & 逐步展开推理增加有效容量 & 代数运算与多步演绎 \\
任务分解 & 串行分治 & $M\times$ 前向 & 递归拆解并逐步求解子目标 & 复杂多步逻辑规划 \\
自一致性 (SC) & 分支采样 & $K\times$ 采样 & 边缘化推导路径取答案众数 & 确定性数学与代码题 \\
思维树/图 (ToT) & 树/图搜索 & $KB\times$ 搜索 & 状态评估、回溯与启发剪枝 & 策略推演与路径规划 \\
验证器选择 & 串联验证 & $(K{+}1)\times$ & 沙箱规则或 PRM 打分选优 & 竞赛代码与定理证明 \\
\bottomrule
\end{tabularx}
\end{table}

\section{推理预算控制}

\subsection{推理总成本}
\term{测试时计算}{Test-Time Compute}{}包括上下文处理、候选生成、验证和工具调用。若每个候选输入长 $L_C$、输出长 $L_i$，不共享前缀时的逻辑词元量为
\begin{equation}
 \begin{aligned}
  N_{\mathrm{in}} &= KL_C, \\
  N_{\mathrm{out}} &= \sum_{i=1}^{K}L_i.
 \end{aligned}
\end{equation}
可用 $C_{\mathrm{cost}}=c_{\mathrm{in}}N_{\mathrm{in}}+c_{\mathrm{out}}N_{\mathrm{out}}+C_v+C_t$ 描述给定口径的成本，其中 $C_v,C_t$ 为验证与工具成本。输入输出单价或等价权重可能不同；词元总数与实际费用之间还差一个价格向量。

共享提示前缀可以降低预填充工作，但每条分支的生成状态仍需维护。并行生成降低墙钟延迟，代价是并发内存上升，总计算量不变。每个请求还需要单独满足上下文约束 $L_C+L_{i,\max}\le L_{\max}$。因此预算至少包含候选上限、输出上限、总计算或费用、时间与工具次数。

增加示例数 $m$ 与增加候选数 $K$ 争用资源。更有用的示例可能提高单候选质量，但较长输入会使每条候选更昂贵；更多采样只有在生成分布包含合适答案且选择器有效时才有意义。\term{边际收益}{Marginal Utility}{}可用开发集上追加一次计算带来的质量改善与成本比较，收益按最终答案质量的改善衡量，并计入生成和选择候选的成本。

\subsection{多数投票的提前停止条件}
若最多还允许生成 $R$ 个候选，当前第一、第二名票数分别为 $n_1,n_2$，且 $n_1>n_2+R$，则剩余候选即使全部投给竞争答案也无法改变唯一第一名。这个条件可以在固定总候选上限下提前结束投票，保持最终最高票决策不变。

该条件固定已有票数与规范化规则，保证最高票决策稳定；答案正确性由任务验证确定。若采用概率置信阈值提前停止，重复查看样本再套用固定样本置信区间会破坏覆盖率；应采用适用于序贯决策的方法，或把停止规则作为完整策略在独立数据上评价。

\begin{algorithm}[有预算的候选聚合与验证]
\AlgoInput{固定上下文、解码策略、解析器、候选与词元上限、最低有效数、平票及验证规则。}
\AlgoOutput{答案、拒答或预算终止状态。}
\AlgoState{答案计数、失败计数、已用词元、验证记录。}
\begin{enumerate}
\item 初始化计数。每次开始生成前，为候选最大输出和可能的验证预留预算；不足时停止追加。
\item 独立采样候选并记录实际停止原因；按固定解析器转为规范答案或失败状态，不静默丢弃失败。
\item 更新票数与预算。若满足预定的验证接受条件或固定上限下不可逆领先条件，进入最终决策。
\item 达到候选上限、超时或预算上限时停止；对平票、有效数不足和未通过验证的状态执行预定升级或拒答策略。
\end{enumerate}
\textbf{不变量：}每个候选都有终态，全部成本被计入，解析与停止规则不根据当前标准答案临时修改。验证接受条件所证明的范围应随结果保留。
\end{algorithm}

\begin{example}


设任务要求计算商品总价：每件 $30$ 元、购买 $6$ 件、优惠 $30$ 元。明确运算为 $30\times6-30=150$ 元。该运算作为候选验证的判定依据。

假设构造上下文含固定指令、问题与模板 $180$ 个词元，两个示例分别为 $90$ 与 $110$ 个词元，因此 $L_C=380$。取 $K=5$，输出长度依次为 $(70,85,60,100,75)$，合计 $390$。按不复用输入、输入输出等权的教学计量，候选成本为 $5\times380+390=2290$ 个词元单位。若一个批量验证调用再耗费 $120$ 个单位，总量为 $2410$。预算按给定输入与输出长度计算。

\begin{table}[htbp]
\centering
\caption[多候选结果与验证状态记录]{多候选结果与验证状态记录。五候选构造示例，仅列最终结果和验证状态。}\label{tab:context-learning-candidates}
\begin{tabularx}{\linewidth}{@{}rllX@{}}
\toprule
候选 & 解析前结果 & 规范答案 & 约束与算术验证\\
\midrule
1 & 150 元 & $150$ 元 & 通过\\
2 & 人民币 150.00 & $150$ 元 & 通过\\
3 & 180 元 & $180$ 元 & 未扣优惠，失败\\
4 & 最终结果缺失 & $\bot$ & 无法验证\\
5 & 120 元 & $120$ 元 & 扣减与题意不符，失败\\
\bottomrule
\end{tabularx}
\end{table}

表\ref{tab:context-learning-candidates}中有效候选为四条，解析覆盖率为 $\frac{4}{5}=80\%$。答案 $150$ 元得两票，整体支持率 $\frac{2}{5}=40\%$，在有效候选中的支持率为 $\frac{2}{4}=50\%$；第二名各一票，票差为一。它是唯一最高票，但不满足严格过半数。若策略预先规定至少四条有效候选且需严格过半数，本次投票应转入验证，事后放宽阈值会破坏规则的预先性。

验证器检查原题运算和单位后，两个等价的 $150$ 元候选通过，其他有效候选失败，可以返回该答案及简洁计算依据。若验证只检查表达式执行，则“$30\times6=180$”也能通过算术执行，仍需检查优惠约束；这说明验证定义直接决定选择结果。

若总预算只有 $2300$ 个单位，当前五候选虽然可以生成，却无法完成预定的 $120$ 单位验证。正确策略应在生成前预留验证成本，或者采用不同候选上限并按其规则决策。验证成本应在生成前计入完整求解预算。
\end{example}

\section{实验迁移性}

\subsection{提示变量隔离}
比较零样本、少样本、分步生成和多候选策略时，应固定模型版本、题目、模板、答案规范化与评分规则。若目标是比较策略在自然成本下的效果，可以允许成本不同但同时报告；若目标是比较预算效率，则应在相同资源上限内比较。两种研究问题都合理，结论只各自适用于对应的问题设定。

使用同一题目作成对比较有助于减少题目难度差异。设两个策略在 $N$ 题上的正确指示为 $c_i^A,c_i^B$，差异估计为
\begin{equation}
\widehat\Delta=\frac1N\sum_i(c_i^A-c_i^B).
\end{equation}
应同时记录 A 对而 B 错、B 对而 A 错的题目，两个总百分比会掩盖这种不对称。采样策略还需要预先规定种子或重复方式；同一问题的多个候选属于同一测试题，据此计算样本量会高估有效样本。

数据应覆盖直接问答、组合运算、约束满足、领域知识与不可回答问题，并按语言、长度和难度报告分子与分母。解析失败、截断、拒答与错误答案是不同状态。只在成功解析的样本上报告准确率，格式成本与覆盖损失就被排除在指标之外。

\subsection{迁移条件校准}
一个提示在某模型上有效，可能依赖其后训练习惯、标签偏好和模板。更换模型、分词器、量化版本、示例顺序或解码参数后，条件分布已经变化。跨模型比较应重新核对这些条件，并采用多题与重复采样评估。

提示开发中反复针对测试集调整示范样本，会使测试分布信息隐式泄漏至运行时配置。评估前须划分独立的开发集并冻结示例库、解析器与终止准则，严防数据污染；近重复样本与同族问题的划分原则参见\bookxref{ch:data-engineering}。此外，数据访问控制与操作授权必须由系统运行时和安全策略网关硬性执行，不能依赖少样本示范或防御性提示词充当安全边界。

\begin{note}[测试时计算的搜索、聚合与验证边界]
分步生成回答“如何组织一条候选”；自一致性回答“哪些答案在采样中出现得更多”；验证选择回答“哪些候选满足指定检查”。三者可以组合，但只有在各自假设、预算和判定范围都明确时，额外计算才具有可解释的收益。
\end{note}


% BEGIN GENERATED INTERVIEWS

% END GENERATED INTERVIEWS
\chapterreferences
\myendchapterdocument
